%%%%%%%%%%%%%%%%%%%%%%%%%%%%%%%%%%%%%%%%%%%%%%%%%%%%%%%%%%%%%%%%%%%%%%%%%%%%%%%%
%% Directional Relational Manifolds, version 6
%% Stratified Variable-Rank Geometry, Rank-Transition Holonomy,
%% Geometric Hysteresis, and Energy-Controlled Transitions
%%%%%%%%%%%%%%%%%%%%%%%%%%%%%%%%%%%%%%%%%%%%%%%%%%%%%%%%%%%%%%%%%%%%%%%%%%%%%%%%

\documentclass[11pt]{article}

\usepackage[utf8]{inputenc}
\usepackage[T1]{fontenc}
\usepackage{amsmath,amssymb,amsthm,mathtools,bm}
\usepackage[margin=1in]{geometry}
\usepackage{enumitem}
\usepackage{booktabs}
\usepackage{array}
\usepackage{xcolor}
\usepackage{microtype}
\usepackage{hyperref}

\newtheorem{definition}{Definition}[section]
\newtheorem{assumption}{Assumption}[section]
\newtheorem{theorem}{Theorem}[section]
\newtheorem{proposition}{Proposition}[section]
\newtheorem{corollary}{Corollary}[section]
\newtheorem{lemma}{Lemma}[section]
\newtheorem{remark}{Remark}[section]
\newtheorem{example}{Example}[section]

\newcommand{\R}{\mathbb{R}}
\newcommand{\E}{\mathcal{E}}
\newcommand{\K}{\mathcal{K}}
\newcommand{\M}{\mathcal{M}}
\newcommand{\Q}{\mathcal{Q}}
\newcommand{\Id}{\operatorname{Id}}
\newcommand{\Ker}{\operatorname{Ker}}
\newcommand{\rank}{\operatorname{rank}}
\newcommand{\Span}{\operatorname{span}}
\newcommand{\im}{\operatorname{im}}
\newcommand{\Hol}{\operatorname{Hol}}
\newcommand{\THol}{\operatorname{THol}}
\newcommand{\End}{\operatorname{End}}
\newcommand{\GL}{\operatorname{GL}}
\newcommand{\Stab}{\operatorname{Stab}}
\newcommand{\dd}{\,\mathrm{d}}
\newcommand{\opnorm}[1]{\left\lVert #1\right\rVert_{\mathrm{op}}}

\hypersetup{
  colorlinks=true,
  linkcolor=blue!65!black,
  citecolor=green!45!black,
  urlcolor=blue!70!black
}

\title{
Directional Relational Manifolds:\\
\large Stratified Variable-Rank Geometry, Rank-Transition Holonomy,\\
Geometric Hysteresis, and Energy-Controlled Transitions
}

\author{
Felipe Maya Muniz\\
Independent Researcher, Brazil\\
\texttt{felipemuniz.grsba@gmail.com}
}

\date{Revised version, July 2026}

\begin{document}
\maketitle

\begin{abstract}
We introduce Directional Relational Manifolds (DRM), a stratified
variable-rank geometric framework for systems whose effective degrees of
freedom may emerge, collapse, or reorganize with state. The construction
starts from a smooth finite-rank vector bundle $E\to M$, a smooth symmetric
positive-semidefinite bundle metric $g$, and an anchor $\rho:E\to TM$. The
effective relational fiber and effective dimension are
\[
\overline E_p=E_p/\Ker g_p,
\qquad
d_{\mathrm{DRM}}(p)=\rank g_p.
\]
The base is decomposed into constant-rank strata. On each stratum, the
effective quotient fibers form a smooth metric vector bundle on which
connections, curvature, parallel transport, energy, and admissible dynamics
are defined. Rank-changing trajectories are modeled as hybrid geometric
paths: invertible metric transport acts inside strata, while possibly
non-invertible transition maps act at rank-changing events.

The global contribution is a theory of \emph{rank-transition holonomy}. For a
closed hybrid path, the ordered composition of intrastatum parallel transports
and interstratum transition maps defines an endomorphism of the initial
effective fiber. Unlike ordinary holonomy, this operator need not be
invertible. We prove that if a closed loop based at an $n$-dimensional
effective fiber contains a transition of rank $r<n$, then the total hybrid
holonomy has rank at most $r$ and therefore cannot be the identity. Thus a
trajectory may return to the same base point and the same nominal dimension
without returning to the same relational state. This produces
\emph{geometric hysteresis}: path-dependent memory generated by dimensional
collapse and reactivation. The effect may occur even when every regular
stratum is flat and every intrastatum transport is trivial.

We further establish an energy criterion for rank transitions. A transition
is dissipative exactly when its metric operator norm is at most one, or,
equivalently in matrix form, $J^{\top}G_+J\preceq G_-$. Compositions of
dissipative transitions and metric parallel transports are globally
non-expansive and generate a contraction semigroup rather than a holonomy
group. Rank defect, norm defect, energy loss, spectral persistence, and
transition commutators are introduced as quantitative hysteresis invariants.
Explicit examples exhibit irreversible rank-$2\to1\to2$ cycles,
order-dependent collapse maps, and memory in flat strata. Under constant-rank,
nondegeneracy, anchor-isomorphism, and compatibility hypotheses, DRM reduces
to Riemannian geometry; with an injective anchor it reduces to
sub-Riemannian geometry; and with the Fisher information metric it reduces to
Fisher--Rao information geometry.
\end{abstract}

\medskip
\noindent\textbf{Keywords:}
directional relational manifolds; variable-rank geometry; degenerate metric;
stratified geometry; rank-transition holonomy; geometric hysteresis;
non-invertible transport; hybrid dynamics; emergent dimensionality;
information geometry.

\medskip
\noindent\textbf{MSC 2020:}
53B05, 53C05, 53C12, 58A30, 37B55, 37C80.

% ==============================================================================
\section{Introduction}

Classical differential geometry normally fixes local dimension before
dynamics is introduced. A smooth $n$-manifold has tangent spaces of constant
dimension $n$, and a Riemannian metric is nondegenerate at every point. These
assumptions are appropriate for many systems, but they do not directly encode
situations in which the number of effective, distinguishable, controllable,
or dynamically available directions changes with state.

Examples include constrained mechanical systems, synchronization and mode
collapse, loss and recovery of controllable directions, rank-deficient
statistical models, changing identifiability, adaptive representations, and
hybrid systems whose internal transition structure reorganizes during
evolution.

Directional Relational Manifolds are proposed as a controlled framework for
such systems. The ambient capacity is represented by a fixed-rank vector
bundle, while the effective geometry is determined by the rank of a
positive-semidefinite relational metric. Directions in the metric kernel are
relationally indistinguishable at the given state. Quotienting by that kernel
produces the effective fiber.

The theory is deliberately stratified. Smooth differential geometry is used
only on constant-rank strata. Rank changes are not forced into a fictitious
globally smooth connection; instead, they are represented by explicit
transition maps. This produces a hybrid geometry consisting of smooth
intrastatum evolution and discrete interstratum reorganization.

The decisive global observation is that a closed base-space loop need not
produce an invertible map of the initial effective fiber. If information is
lost during an intermediate dimensional collapse, deterministic later
expansion cannot reconstruct that information from the surviving state alone.
The resulting closed-loop transport is therefore naturally semigroup-valued,
not group-valued.

\subsection{Main contributions}

No novelty is claimed for positive-semidefinite metrics, rank
stratifications, quotient bundles, anchors, connections, rank inequalities,
or hybrid reset maps considered separately. The contribution is their
integration into the chain
\[
\begin{aligned}
\text{variable metric rank}
&\longrightarrow \text{effective relational dimension}
\\
&\longrightarrow \text{stratumwise quotient geometry}
\\
&\longrightarrow \text{hybrid rank-transition transport}
\\
&\longrightarrow \text{non-invertible holonomy and hysteresis}.
\end{aligned}
\]

The principal results are:
\begin{enumerate}[label=(\roman*)]
\item a precise stratified DRM structure based on effective quotient fibers;
\item an ordered hybrid transport law for paths visiting multiple rank strata;
\item a transition-holonomy semigroup for closed hybrid loops;
\item a dimensional holonomy theorem forcing rank-deficient loop transport
whenever a genuine rank-reducing transition occurs;
\item a proof that geometric hysteresis can exist with zero intrastatum
curvature;
\item energy and matrix criteria for dissipative, isometric, and expansive
rank transitions;
\item quantitative memory invariants based on rank, norm, energy, spectrum,
and path order;
\item reduction theorems connecting DRM to established geometries.
\end{enumerate}

\subsection{Organization}

Section~\ref{sec:related} positions DRM relative to existing frameworks.
Sections~\ref{sec:ambient}--\ref{sec:strata} define the ambient bundle,
relational metric, quotient fibers, and rank stratification.
Sections~\ref{sec:anchor}--\ref{sec:connection} introduce anchored motion,
connections, curvature, and regular transport.
Sections~\ref{sec:hybrid}--\ref{sec:holonomy} define hybrid paths and
rank-transition holonomy.
Sections~\ref{sec:hysteresis}--\ref{sec:order} develop geometric hysteresis,
energy control, and order dependence.
Section~\ref{sec:geodesics} introduces relational action.
Sections~\ref{sec:examples} and~\ref{sec:reductions} give examples and
reduction theorems.

% ==============================================================================
\section{Relation to Existing Mathematical Frameworks}
\label{sec:related}

\subsection{Degenerate metrics}

A smooth metric may become degenerate, preventing use of the ordinary inverse
metric and obstructing standard Levi--Civita constructions. Singular
semi-Riemannian geometry develops curvature-like objects under additional
regularity hypotheses \cite{Stoica2014}. DRM takes a different route: on each
constant-rank stratum it quotients the ambient fiber by the metric kernel,
thereby obtaining a positive-definite effective metric. It does not claim to
replace intrinsic theories of degenerate metrics; it supplies a
quotient-and-transition framework adapted to changing effective rank.

\subsection{Stratified spaces, singular distributions, and holonomy groupoids}

Stratified geometry studies spaces decomposed into smooth pieces of different
dimensions. Singular distributions and singular foliations permit pointwise
variation in rank and possess sophisticated orbit and holonomy theories
\cite{Sussmann1973,AndroulidakisSkandalis2009}. DRM instead stratifies a smooth
ambient base by the rank of a semidefinite bundle metric and assigns an
effective quotient bundle to each stratum. Its interstratum maps are explicit
model data and may be non-invertible. Consequently, the global transport
object is generally a category or semigroup of linear maps, rather than a
groupoid of locally invertible arrows.

\subsection{Anchored bundles and Lie algebroids}

An anchored vector bundle consists of a vector bundle and a bundle map into
the tangent bundle. A Lie algebroid adds a compatible bracket on sections
\cite{Mackenzie2005}. DRM uses anchors to convert internal relational
directions into observable state velocities. A Lie-algebroid structure may be
imposed on a regular stratum when intrinsic torsion and curvature are needed,
but it is not part of the minimal global definition.

\subsection{Sub-Riemannian geometry}

Sub-Riemannian geometry equips a distribution of admissible directions with a
positive-definite metric. On a constant-rank stratum with injective effective
anchor, a DRM induces a sub-Riemannian structure. DRM additionally permits the
effective distribution rank to change through metric degeneration and
incorporates explicit reset maps between rank strata.

\subsection{Information geometry}

Information geometry equips regular statistical models with the Fisher metric
\cite{AmariNagaoka2000,Chentsov1982,Nielsen2020}. Rank-deficient Fisher matrices occur in
non-identifiable or singular models. DRM interprets the metric rank as an
effective relational dimension and treats passages through singular rank loci
using stratification and transition maps. Fisher--Rao geometry is recovered in
the regular constant-rank case.

\subsection{Hybrid dynamical systems}

Hybrid systems combine continuous flows with discrete jumps or reset maps
\cite{GoebelSanfeliceTeel2012}. DRM uses this architecture geometrically:
connections and flows act inside rank strata, while transition maps act on
effective fibers at rank-changing events. The maps are assessed not only
dynamically but also metrically, through energy inequalities.

\subsection{Classical holonomy}

For a connection on a fixed-rank vector bundle, parallel transport around a
closed loop gives an invertible endomorphism, and the holonomy at a base point
forms a subgroup of $\GL(E_p)$ \cite{KobayashiNomizu1963,Lee2018}. In DRM,
regular transport remains invertible inside each stratum, but interstratum
maps may lower rank. The ordered closed-loop transport can therefore lie in
$\End(\overline E_p)\setminus\GL(\overline E_p)$. The qualifier
``rank-transition'' distinguishes this object from ordinary holonomy.

% ==============================================================================
\section{Ambient Relational Structure}
\label{sec:ambient}

\begin{definition}[Ambient state manifold]
Let $M$ be a smooth, Hausdorff, second-countable manifold.
\end{definition}

\begin{definition}[Ambient relational bundle]
Let $\pi:E\to M$ be a smooth real vector bundle of finite constant rank $N$.
\end{definition}

The rank $N$ is the maximum ambient relational capacity. It is not the
effective local dimension.

\begin{definition}[Relational metric]
A relational metric is a smooth section
\[
g\in\Gamma(\operatorname{Sym}^2E^*)
\]
such that every bilinear form $g_p:E_p\times E_p\to\R$ is symmetric and
positive semidefinite.
\end{definition}

\begin{definition}[Relational kernel]
At $p\in M$, define
\[
K_p:=\Ker g_p
=\{v\in E_p:g_p(v,w)=0\text{ for every }w\in E_p\}.
\]
\end{definition}

\begin{definition}[Effective relational fiber]
The effective fiber at $p$ is
\[
\overline E_p:=E_p/K_p.
\]
\end{definition}

\begin{proposition}[Induced effective metric]
The formula
\[
\overline g_p([v],[w]):=g_p(v,w)
\]
defines a positive-definite inner product on $\overline E_p$.
\end{proposition}

\begin{proof}
If $v'=v+k_1$ and $w'=w+k_2$ with $k_1,k_2\in K_p$, every term containing
$k_1$ or $k_2$ vanishes, so the value is independent of representatives. If
$\overline g_p([v],[v])=0$, the Cauchy--Schwarz inequality for positive
semidefinite forms gives $g_p(v,w)=0$ for every $w$, hence $v\in K_p$ and
$[v]=0$.
\end{proof}

\begin{definition}[Effective relational dimension]
The DRM dimension at $p$ is
\[
d_{\mathrm{DRM}}(p):=\rank g_p=\dim\overline E_p.
\]
\end{definition}

\begin{remark}
Dimension is not the cardinality of a directional generating set. Rank is the
invariant measuring the number of metrically distinguishable directions.
\end{remark}

% ==============================================================================
\section{Rank Stratification}
\label{sec:strata}

For $r\in\{0,\ldots,N\}$, define
\[
M_r:=\{p\in M:\rank g_p=r\}.
\]

\begin{assumption}[Regular rank stratification]
Each connected component of every nonempty $M_r$ is a smooth embedded
submanifold of $M$, and the collection of components is locally finite.
\end{assumption}

A connected component $S\subseteq M_r$ is called a rank-$r$ stratum.

\begin{proposition}[Kernel subbundle]
On a rank-$r$ stratum $S$, the union
\[
K|_S:=\bigcup_{p\in S}K_p
\]
is a smooth vector subbundle of $E|_S$ of rank $N-r$.
\end{proposition}

\begin{proof}
The metric induces a smooth bundle morphism
\[
g^\flat:E|_S\to E^*|_S,
\qquad g^\flat_p(v)=g_p(v,\cdot).
\]
Its rank is constantly $r$ on $S$. The constant-rank theorem for vector-bundle
morphisms \cite{Lee2013} implies that $\Ker g^\flat=K|_S$ is a smooth
subbundle.
\end{proof}

\begin{corollary}[Effective quotient bundle]
The quotient
\[
\overline E^{(S)}:=E|_S/K|_S
\]
is a smooth rank-$r$ vector bundle over $S$, equipped with the smooth
positive-definite metric $\overline g^{(S)}$.
\end{corollary}

The quotient bundles generally do not assemble into a single vector bundle
over all of $M$, because their ranks differ.

% ==============================================================================
\section{Anchor and Admissible Motion}
\label{sec:anchor}

\begin{definition}[Anchor]
An anchor is a smooth vector-bundle morphism $\rho:E\to TM$ covering the
identity on $M$.
\end{definition}

\begin{assumption}[Kernel compatibility and stratum tangency]
For every $p\in M$,
\[
K_p\subseteq\Ker\rho_p.
\]
Moreover, for every rank stratum $S$ and every $p\in S$,
\[
\rho_p(E_p)\subseteq T_pS\subseteq T_pM.
\]
\end{assumption}

\begin{proposition}[Effective stratum anchor]
Kernel compatibility defines a unique map
\[
\overline\rho_p:\overline E_p\to T_pM,
\qquad \overline\rho_p([v])=\rho_p(v).
\]
On every stratum $S$, the tangency assumption makes these maps a smooth
bundle morphism
\[
\overline\rho^{(S)}:\overline E^{(S)}\to TS.
\]
Consequently, $\overline\rho(U)[f]$ is intrinsically defined for
$U\in\Gamma(\overline E^{(S)})$ and $f\in C^\infty(S)$.
\end{proposition}

\begin{definition}[Admissible curve]
Let $S$ be a stratum. A $C^1$ curve $\gamma:I\to S$ is admissible if there
exists a smooth field $V(t)\in\overline E^{(S)}_{\gamma(t)}$ such that
\[
\dot\gamma(t)=\overline\rho_{\gamma(t)}V(t).
\]
\end{definition}

The effective anchor may be neither injective nor surjective. Internal
directions can be unobservable, and tangent directions can be inaccessible.

% ==============================================================================
\section{Stratumwise Connections, Curvature, and Regular Transport}
\label{sec:connection}

Let $S$ be a fixed rank stratum.

\begin{assumption}[Anchored bracket]
The bundle $\overline E^{(S)}$ is equipped with a bracket
\[
[\cdot,\cdot]_S:\Gamma(\overline E^{(S)})\times
\Gamma(\overline E^{(S)})\to\Gamma(\overline E^{(S)})
\]
satisfying
\[
[U,fV]_S=f[U,V]_S+\overline\rho(U)[f]V,
\qquad
\overline\rho([U,V]_S)=[\overline\rho(U),\overline\rho(V)].
\]
\end{assumption}

If the Jacobi identity also holds, the stratum carries a Lie-algebroid
structure.

\begin{definition}[Relational connection]
A relational connection on $\overline E^{(S)}$ is a map
\[
\nabla^{(S)}:\Gamma(\overline E^{(S)})\times
\Gamma(\overline E^{(S)})\to\Gamma(\overline E^{(S)})
\]
that is $C^\infty(S)$-linear in the first argument and satisfies
\[
\nabla^{(S)}_U(fV)=\overline\rho(U)[f]V+f\nabla^{(S)}_UV.
\]
\end{definition}

\begin{definition}[Metric compatibility]
The connection is metric-compatible when
\[
\overline\rho(U)[\overline g(V,W)]
=\overline g(\nabla^{(S)}_UV,W)
+\overline g(V,\nabla^{(S)}_UW).
\]
\end{definition}

\begin{definition}[Torsion and curvature]
Define
\[
T^{(S)}(U,V)=\nabla^{(S)}_UV-\nabla^{(S)}_VU-[U,V]_S
\]
and
\[
R^{(S)}(U,V)W
=\nabla^{(S)}_U\nabla^{(S)}_VW
-\nabla^{(S)}_V\nabla^{(S)}_UW
-\nabla^{(S)}_{[U,V]_S}W.
\]
\end{definition}

\begin{proposition}[Existence of metric-compatible connections]
Every effective metric bundle $(\overline E^{(S)},\overline g^{(S)})$ admits a
metric-compatible connection.
\end{proposition}

\begin{proof}
Choose any connection $\nabla^0$ and define
\[
Q(U;V,W)=\overline\rho(U)[\overline g(V,W)]
-\overline g(\nabla^0_UV,W)-\overline g(V,\nabla^0_UW).
\]
Because $\overline g$ is nondegenerate, there is a unique endomorphism-valued
one-form $A$ satisfying
\[
2\overline g(A_UV,W)=Q(U;V,W).
\]
Then $\nabla_UV:=\nabla^0_UV+A_UV$ is metric-compatible.
\end{proof}

A metric-compatible connection defines parallel transport
\[
P^\gamma_{s\to t}:\overline E^{(S)}_{\gamma(s)}
\to\overline E^{(S)}_{\gamma(t)}
\]
along admissible curves in $S$. It is an isometry and therefore invertible.
For a closed loop $\ell$ in $S$ based at $p$,
\[
\Hol^{(S)}_\ell:=P^\ell_{0\to1}\in\GL(\overline E_p)
\]
is the ordinary stratumwise holonomy.

% ==============================================================================
\section{Hybrid Paths and Ordered Transport}
\label{sec:hybrid}

\begin{definition}[Rank-transition event]
Let $p_-\in S_-$ and $p_+\in S_+$ lie in strata of ranks $r$ and $s$,
respectively. A rank-transition event consists of the ordered base reset
$p_-\rightsquigarrow p_+$ together with a linear map
\[
J_{r\to s}:\overline E_{p_-}\to\overline E_{p_+}.
\]
The points $p_-$ and $p_+$ need not coincide, and the map need not be
invertible. Thus an event records both the reset of the base state and the
reset of the effective relational state.
\end{definition}

\begin{definition}[Hybrid DRM path]
A hybrid DRM path is a finite sequence
\[
\gamma=(\gamma_1,J_1,\gamma_2,J_2,\ldots,J_{m-1},\gamma_m),
\]
where each $\gamma_i:[a_i,b_i]\to S_i$ is a smooth admissible path contained
in one rank stratum and each $J_i$ is a transition event
\[
J_i:\overline E_{\gamma_i(b_i)}
\longrightarrow\overline E_{\gamma_{i+1}(a_{i+1})}.
\]
No continuity of the base path is imposed at a transition event. A continuous
crossing of a rank boundary can instead be treated only after specifying
one-sided trace fibers and limiting parallel transports; those additional
analytic data are not assumed in the present finite-event framework.
\end{definition}

Let $P_i$ denote parallel transport along $\gamma_i$.

\begin{definition}[Ordered hybrid transport]
The hybrid transport operator is
\[
\mathcal P_\gamma
=P_mJ_{m-1}P_{m-1}\cdots J_2P_2J_1P_1.
\]
\end{definition}

The order is essential. The maps may be noncommuting and may act between
spaces of different dimensions.

\begin{proposition}[Concatenation law]
If $\gamma$ ends where $\eta$ begins, then
\[
\mathcal P_{\eta*\gamma}=\mathcal P_\eta\mathcal P_\gamma.
\]
\end{proposition}

\begin{proof}
A vector is first transported along $\gamma$ and then along $\eta$. This is
exactly composition of the ordered operators.
\end{proof}

\begin{remark}[Path category]
If rank strata are represented as vertices and permitted transitions as
directed edges, hybrid paths form a path category. Sending each vertex to its
effective fiber and each path to its ordered transport defines a functor into
the category of finite-dimensional vector spaces and linear maps. The target
is not generally a groupoid because transition maps need not be invertible.
\end{remark}

% ==============================================================================
\section{Rank-Transition Holonomy}
\label{sec:holonomy}

\begin{definition}[Closed hybrid loop]
A hybrid path $\gamma$ is closed and based at $p$ if its initial and final
base points coincide with $p$, and the initial and final effective fibers are
identified with $\overline E_p$.
\end{definition}

\begin{definition}[Rank-transition holonomy]
For a closed hybrid loop $\gamma$ based at $p$, define
\[
H_\gamma:=\mathcal P_\gamma\in\End(\overline E_p).
\]
Unlike ordinary holonomy, $H_\gamma$ need not belong to
$\GL(\overline E_p)$.
\end{definition}

\begin{definition}[Transition-holonomy semigroup]
The transition-holonomy semigroup at $p$ is
\[
\THol_p(\mathfrak D)
:=\{H_\gamma:\gamma\text{ is a closed hybrid DRM loop based at }p\}
\subseteq\End(\overline E_p).
\]
\end{definition}

\begin{proposition}
$\THol_p(\mathfrak D)$ contains the identity and is closed under composition.
It need not contain inverses.
\end{proposition}

\begin{proof}
The constant loop gives the identity. Concatenation of closed loops gives
composition by the concatenation law. If a loop holonomy has deficient rank,
it is not invertible and no inverse belongs to $\End(\overline E_p)$ as an
inverse map.
\end{proof}

\begin{theorem}[Dimensional holonomy theorem]
\label{thm:rank-holonomy}
Let $\gamma$ be a closed hybrid DRM loop based at $p$, with
$\dim\overline E_p=n$. Suppose one transition map $J_k$ appearing in the
ordered composition has rank $r<n$. Then
\[
\rank H_\gamma\le r<n.
\]
Consequently,
\[
H_\gamma\ne\Id_{\overline E_p},
\]
and $H_\gamma$ is non-invertible.
\end{theorem}

\begin{proof}
Write the ordered composition as
\[
H_\gamma=A J_k B,
\]
where $B$ contains every map applied before $J_k$ and $A$ contains every map
applied after it. The rank inequality for compositions gives
\[
\rank(AJ_kB)\le\rank(J_k)=r.
\]
The identity on an $n$-dimensional space has rank $n$, so $H_\gamma$ cannot be
the identity and cannot be invertible.
\end{proof}

\begin{corollary}[Irreversible dimensional cycle]
A closed loop containing a genuine rank-reducing transition cannot recover the
identity transport by later deterministic linear maps acting only on the
surviving effective state.
\end{corollary}

\begin{proof}
Once the rank of the total map has fallen below $n$, postcomposition cannot
increase it above the rank bound established in
Theorem~\ref{thm:rank-holonomy}.
\end{proof}

\begin{remark}[External memory]
Recovery is possible only after enlarging the state with an auxiliary channel
that stores the discarded information. The enlarged evolution may be
invertible even though the projected effective evolution is not.
\end{remark}

% ==============================================================================
\section{Geometric Hysteresis and Memory Invariants}
\label{sec:hysteresis}

\begin{definition}[Geometric hysteresis]
A closed hybrid loop exhibits geometric hysteresis if
\[
H_\gamma\ne\Id.
\]
It exhibits irreversible geometric hysteresis if $H_\gamma$ is non-invertible.
\end{definition}

A system exhibiting hysteresis returns to the same base point but not to the
same effective relational state.

\begin{definition}[Surviving and lost subspaces]
For a closed loop $\gamma$, define
\[
\mathcal M_\gamma:=\im H_\gamma,
\qquad
\mathcal L_\gamma:=\Ker H_\gamma.
\]
\end{definition}

\begin{definition}[Rank-memory defect]
If $\dim\overline E_p=n$, define
\[
\delta_{\mathrm{rank}}(\gamma):=n-\rank H_\gamma.
\]
\end{definition}

\begin{definition}[Norm-memory defect]
Define
\[
\delta_{\mathrm{op}}(\gamma):=
\opnorm{H_\gamma-\Id}.
\]
A normalized Frobenius alternative is
\[
h_F(\gamma):=
\frac{\|H_\gamma-\Id\|_F}{\sqrt n}.
\]
\end{definition}

\begin{theorem}[Hysteresis without intrastatum curvature]
\label{thm:flat-hysteresis}
Let $\gamma$ be a closed hybrid DRM loop. Assume:
\begin{enumerate}[label=(\roman*)]
\item every regular stratum visited by $\gamma$ is flat along its regular segments;
\item every intrastatum transport is the identity;
\item at least one transition map has rank below the initial fiber dimension
$n$.
\end{enumerate}
Then $H_\gamma$ is non-invertible and different from the identity.
\end{theorem}

\begin{proof}
The identity transports drop out, leaving an ordered product of transition
maps. Theorem~\ref{thm:rank-holonomy} applies.
\end{proof}

\begin{remark}
This separates two sources of path dependence:
\[
\text{intrastatum curvature}
\qquad\text{and}\qquad
\text{interstratum dimensional memory}.
\]
The second can remain nontrivial when the first vanishes.
\end{remark}

% ==============================================================================
\section{Energy-Controlled Rank Transitions}
\label{sec:energy}

Let the incoming and outgoing effective fibers be finite-dimensional inner
product spaces
\[
(V_-,\langle\cdot,\cdot\rangle_-),
\qquad
(V_+,\langle\cdot,\cdot\rangle_+).
\]

\begin{definition}[Transition energy]
For $v\in V_-$, define
\[
\mathcal E_-(v)=\frac12\|v\|_-^2,
\qquad
\mathcal E_+(Jv)=\frac12\|Jv\|_+^2.
\]
\end{definition}

\begin{definition}[Transition classes]
A transition $J:V_-\to V_+$ is:
\begin{enumerate}[label=(\roman*)]
\item \emph{dissipative} if $\|Jv\|_+\le\|v\|_-$ for every $v$;
\item \emph{isometric} if equality holds for every $v$;
\item \emph{expansive} if the dissipative inequality fails for some $v$.
\end{enumerate}
\end{definition}

Define the metric operator norm
\[
\|J\|_{-\to+}:=
\sup_{v\ne0}\frac{\|Jv\|_+}{\|v\|_-}.
\]

\begin{proposition}[Transition-energy criterion]
A transition $J$ is dissipative if and only if
\[
\|J\|_{-\to+}\le1.
\]
It is isometric if and only if
\[
J^*\overline g_+=\overline g_-.
\]
\end{proposition}

\begin{proof}
The first assertion is the definition of the induced operator norm. Equality
of norms for every vector is equivalent, by polarization, to preservation of
the full inner product, which is $J^*\overline g_+=\overline g_-$.
\end{proof}

\begin{corollary}[Rank-loss obstruction to isometry]
If $\dim V_+<\dim V_-$, no map $J:V_-\to V_+$ is isometric on all of $V_-$.
\end{corollary}

\begin{corollary}[Canonical dissipative collapse]
Let $W\subseteq V_-$ and let $\Pi_W:V_-\to W$ be the orthogonal projection.
If $V_+=W$ with the restricted metric, then $\Pi_W$ is dissipative and
\[
\|v\|_-^2=\|\Pi_Wv\|_+^2+\|v-\Pi_Wv\|_-^2.
\]
\end{corollary}

\begin{proposition}[Compositional stability]
A composition of dissipative transitions is dissipative.
\end{proposition}

\begin{proof}
The induced operator norm is submultiplicative.
\end{proof}

\begin{theorem}[Global non-expansiveness]
Consider a hybrid path composed of metric-compatible parallel transports
$P_i$ and dissipative transitions $J_i$. Then
\[
\|\mathcal P_\gamma v\|_{\mathrm{final}}
\le\|v\|_{\mathrm{initial}}
\]
for every initial vector $v$.
\end{theorem}

\begin{proof}
Every metric parallel transport has norm one and every dissipative transition
has norm at most one. Apply submultiplicativity to the ordered product.
\end{proof}

\begin{corollary}[Contraction semigroup]
Dissipative closed-loop holonomies form a contraction semigroup in
$\End(\overline E_p)$.
\end{corollary}

\begin{definition}[Energy hysteresis]
For a closed loop based at $p$, define
\[
\Delta_\gamma(v)
:=\frac12\|v\|^2-\frac12\|H_\gamma v\|^2.
\]
For dissipative loops, $\Delta_\gamma(v)\ge0$.
\end{definition}

In a basis with incoming metric matrix $G_-$ and outgoing metric matrix $G_+$,
dissipativity is equivalent to
\[
J^\top G_+J\preceq G_-.
\]
For a closed loop with metric matrix $G$, define the energy-defect matrix
\[
D_\gamma:=G-H_\gamma^\top G H_\gamma.
\]
The loop is dissipative exactly when $D_\gamma\succeq0$.

\begin{remark}[Dimension activation]
When $\dim V_+>\dim V_-$, a transition may be isometric on its image, but the
new complementary directions are not determined by the incoming vector.
Forcing, noise, initialization data, a learned rule, or an auxiliary memory
channel is needed to populate them.
\end{remark}

% ==============================================================================
\section{Path-Order Dependence and Noncommutative Memory}
\label{sec:order}

\begin{definition}[Transition commutator]
For endomorphisms $A,B$ of a common effective fiber, define
\[
[A,B]_{\mathrm{tr}}:=AB-BA.
\]
\end{definition}

\begin{definition}[Order-memory defect]
Define
\[
\delta_{\mathrm{ord}}(A,B):=\|AB-BA\|,
\]
and the normalized index
\[
h_{\mathrm{ord}}(A,B)
:=\frac{\|AB-BA\|_F}
{\max\{1,\|A\|_F\|B\|_F\}}.
\]
\end{definition}

\begin{proposition}[Noncommutative dimensional memory]
If $[A,B]_{\mathrm{tr}}\ne0$, then there exists an initial vector for which
the ordered histories $AB$ and $BA$ produce different final states.
\end{proposition}

\begin{proof}
A nonzero linear map $AB-BA$ acts nontrivially on at least one vector.
\end{proof}

Thus the same local transition operations can encode different memory when
traversed in a different order.

% ==============================================================================
\section{Spectral and Asymptotic Memory}
\label{sec:spectral}

Repeated traversal of a closed loop evolves the state by
\[
v_k=H_\gamma^k v_0.
\]

\begin{proposition}[Strict contraction]
If $\opnorm{H_\gamma}<1$, then $H_\gamma^k v\to0$ for every $v$.
\end{proposition}

\begin{proof}
$\|H_\gamma^k v\|\le\opnorm{H_\gamma}^k\|v\|\to0$.
\end{proof}

\begin{proposition}[Projection memory]
If $H_\gamma^2=H_\gamma$, then
\[
H_\gamma^k v=H_\gamma v
\qquad(k\ge1).
\]
The persistent memory space is $\im H_\gamma$.
\end{proposition}

More generally, the spectrum and Jordan structure of $H_\gamma$ determine
asymptotic persistence, decay, and possible transient amplification.

% ==============================================================================
\section{Relational Energy, Length, and Hybrid Action}
\label{sec:geodesics}

Let $\gamma:[a,b]\to S$ be admissible in one stratum and satisfy
$\dot\gamma=\overline\rho(V)$.

\begin{definition}[Relational energy and length]
Define
\[
\mathcal E[\gamma,V]
=\frac12\int_a^b\overline g(V,V)\dd t,
\qquad
L[\gamma,V]
=\int_a^b\sqrt{\overline g(V,V)}\dd t.
\]
\end{definition}

\begin{definition}[Stratumwise relational geodesic]
An admissible pair $(\gamma,V)$ is a stratumwise relational geodesic if it is
stationary for $\mathcal E$ under admissible fixed-endpoint variations.
\end{definition}

For a hybrid path with segments $(\gamma_i,V_i)$ and transitions $J_i$, define
\[
\mathcal A_{\mathrm{hyb}}
=\sum_i\mathcal E[\gamma_i,V_i]
+\sum_i C_i(v_i^-,v_i^+),
\]
where a natural cost for a dissipative transition is
\[
C_i
=\frac12\|v_i^-\|_-^2-
\frac12\|J_iv_i^-\|_+^2\ge0.
\]
This separates continuous action inside strata from energy lost at discrete
rank changes.

% ==============================================================================
\section{Examples}
\label{sec:examples}

\subsection{A corrected rank-$2\to1\to2$ cycle}

Let the initial and final effective fibers be $\R^2$. Define
\[
J_{2\to1}(a,b)=a,
\qquad
J_{1\to2}(a)=(a,0),
\]
and assume all regular transports are identities. The closed-loop holonomy is
\[
H=J_{1\to2}J_{2\to1}
=\begin{pmatrix}1&0\\0&0\end{pmatrix}.
\]
Therefore
\[
H(a,b)=(a,0),
\qquad
\rank H=1,
\qquad
\delta_{\mathrm{rank}}=1.
\]
The energy loss is
\[
\Delta(a,b)=\frac12(a^2+b^2)-\frac12a^2=\frac12b^2.
\]
The map is idempotent, so one traversal projects the state onto a persistent
one-dimensional memory subspace.

\subsection{A smooth metric realizing a rank boundary}

Let
\[
M=\R,
\qquad
E=M\times\R^2,
\qquad
g_x=\begin{pmatrix}1&0\\0&x^2\end{pmatrix}.
\]
Then $\rank g_x=2$ for $x\ne0$ and $\rank g_0=1$. The rank-two strata are
$(-\infty,0)$ and $(0,\infty)$, and the rank-one stratum is $\{0\}$. At the
singular point,
\[
K_0=\Span\{e_2\},
\qquad
\overline E_0\cong\R[e_1].
\]
Define
\[
\rho_x(ae_1+be_2)=xa\,\partial_x.
\]
For $x\ne0$ the image lies in the tangent bundle of the corresponding open
rank-two stratum. At $x=0$, $\rho_0=0$, so both
$K_0\subseteq\Ker\rho_0$ and
$\rho_0(E_0)\subseteq T_0\{0\}=\{0\}$ hold.

In the finite-event hybrid model, choose $x_-<0$ and $x_+>0$. The passage is
represented by a reset from $x_-$ to $0$ with
$J_{2\to1}(a,b)=a$, followed by a reset from $0$ to $x_+$ with
$J_{1\to2}(a)=(a,0)$. This is a hybrid idealization of passage through the
rank boundary, not a claim that a compact path lying in one open stratum has
$0$ as an endpoint. The second component is discarded at entry and is not
reconstructed at exit.

\subsection{Two noncommuting collapse directions}

Let
\[
\Pi_x=\begin{pmatrix}1&0\\0&0\end{pmatrix},
\qquad
u_\theta=\begin{pmatrix}\cos\theta\\\sin\theta\end{pmatrix},
\qquad
\Pi_\theta=u_\theta u_\theta^\top.
\]
For generic $\theta$,
\[
\Pi_x\Pi_\theta
=\begin{pmatrix}
\cos^2\theta&\cos\theta\sin\theta\\
0&0
\end{pmatrix},
\]
whereas
\[
\Pi_\theta\Pi_x
=\begin{pmatrix}
\cos^2\theta&0\\
\cos\theta\sin\theta&0
\end{pmatrix}.
\]
Thus $\Pi_x\Pi_\theta\ne\Pi_\theta\Pi_x$, and the final state depends on the
order of collapse operations.

\subsection{Flat strata with nontrivial hysteresis}

Let the base be a closed planar path divided into flat strata. Give every
stratum a trivial connection, so every $P_i=\Id$. Insert the collapse and
re-embedding maps of the rank-$2\to1\to2$ example. Then every regular curvature
vanishes, yet
\[
H=\begin{pmatrix}1&0\\0&0\end{pmatrix}\ne\Id.
\]
The memory is generated exclusively by rank transitions.

\subsection{External memory restores the discarded component}

Enlarge the state by an auxiliary variable $\xi$ and define
\[
\widetilde J_{2\to1}(a,b)=(a,\xi=b).
\]
A later activation map may use the stored value:
\[
\widetilde J_{1\to2}(a,\xi)=(a,\xi).
\]
The enlarged evolution recovers $(a,b)$, but recovery does not occur inside
the one-dimensional effective fiber alone. This identifies precisely where
additional information must enter.

\subsection{A non-flat regular DRM}

Let
\[
M=\R^2,
\qquad
E=TM,
\qquad
\rho=\Id,
\]
and set
\[
\phi(x,y)=x^2+y^2,
\qquad
g=e^{2\phi}(dx^2+dy^2).
\]
There is one rank-two stratum. The DRM connection is the Levi--Civita
connection and the Gaussian curvature is
\[
K=-e^{-2\phi}\Delta\phi
=-4e^{-2(x^2+y^2)}.
\]
This verifies that a regular DRM reproduces ordinary curved Riemannian
geometry.

\subsection{Energy-amplifying activation}

Let $V_-=\R$ and $V_+=\R^2$ with Euclidean metrics. Define
\[
J_\alpha(a)=(a,\alpha a).
\]
Then
\[
\|J_\alpha a\|_+^2=(1+\alpha^2)a^2.
\]
For $\alpha\ne0$, the transition is expansive. The additional energy must be
supplied by an external source or represented by an activation cost.

% ==============================================================================
\section{Reduction Theorems}
\label{sec:reductions}

\begin{theorem}[Riemannian reduction]
Suppose:
\begin{enumerate}[label=(\roman*)]
\item $M$ is an $n$-manifold;
\item $E\to M$ has rank $n$;
\item $g$ is positive definite everywhere;
\item $\rho:E\to TM$ is a vector-bundle isomorphism;
\item the DRM connection is metric-compatible and torsion-free.
\end{enumerate}
Then DRM induces a Riemannian metric $h$ on $M$, and the transferred DRM
connection is the Levi--Civita connection of $h$.
\end{theorem}

\begin{proof}
Since $g$ is positive definite, $K=0$ and $\overline E=E$. Define
\[
h_p(X,Y)=g_p(\rho_p^{-1}X,\rho_p^{-1}Y).
\]
This is a smooth Riemannian metric. Transfer the connection through $\rho$.
The resulting connection on $TM$ is metric-compatible and torsion-free;
uniqueness of the Levi--Civita connection gives the result.
\end{proof}

\begin{corollary}[Sub-Riemannian reduction]
On a constant-rank stratum, if the effective anchor is injective, its image
\[
\mathcal H:=\im\overline\rho\subseteq TS
\]
is a constant-rank distribution. Transporting $\overline g$ through the
anchor defines a sub-Riemannian metric on $\mathcal H$.
\end{corollary}

\begin{corollary}[Fisher--Rao reduction]
Let $M=\Theta$ be a regular $n$-dimensional parameter manifold. Assume the
effective anchor identifies $\overline E_\theta$ with $T_\theta\Theta$ and
\[
\overline g_\theta(e_i,e_j)
=\mathbb E_\theta\!\left[
\partial_i\log p(X;\theta)\,
\partial_j\log p(X;\theta)
\right].
\]
If the connection is metric-compatible and torsion-free, the effective DRM is
the Fisher--Rao statistical manifold.
\end{corollary}

% ==============================================================================
\section{Interpretation and Potential Applications}

DRM separates three quantities that are often conflated:
\begin{enumerate}
\item \emph{ambient capacity}: the bundle rank $N$;
\item \emph{effective dimension}: the metric rank $d_{\mathrm{DRM}}(p)$;
\item \emph{observable mobility}: the rank of the effective anchor
$\overline\rho_p$.
\end{enumerate}

These need not agree. A system may retain a large ambient state capacity while
its metric rank collapses. It may retain high metric rank while only a few
directions are dynamically observable.

Potential interpretations include:
\begin{itemize}
\item temporary loss of a mechanical or controllable degree of freedom;
\item passage through a non-identifiable statistical stratum;
\item collapse and re-expansion of latent representations;
\item switching systems in which the order of constraints matters;
\item memory induced by dimensional reduction rather than ordinary curvature.
\end{itemize}

The mathematical framework does not by itself establish any specific physical,
cognitive, or computational application. Such claims require an empirical
identification procedure for effective fibers, rank boundaries, transition
maps, and memory defects.

% ==============================================================================
\section{Limitations and Open Problems}
\label{sec:limitations}

\paragraph{Regularity of rank loci.}
The regular-stratification assumption excludes pathological rank sets.
Whitney stratifications, differential spaces, diffeology, or sheaf-theoretic
methods may be required in greater generality.

\paragraph{Canonical transition maps.}
Transition maps are prescribed data. A theory deriving them canonically from
local metric degeneration, dynamics, or a variational principle remains open.

\paragraph{Continuous rank-boundary crossings.}
The finite-event model uses explicit base resets between points in different
strata. A genuinely continuous crossing requires one-sided trace fibers,
limiting parallel transports, and hypotheses ensuring that those limits exist
and are independent of auxiliary trivializations. That extension is not
developed here.

\paragraph{Relation to singular foliation holonomy.}
A detailed comparison with holonomy groupoids of singular foliations is still
needed. DRM explicitly admits non-invertible interstratum maps, but the exact
categorical relationship deserves further study.

\paragraph{Nonlinear transitions.}
The present theory is linear at the effective-fiber level. Nonlinear reset
maps require differential, tangent, or categorical extensions.

\paragraph{Infinite transition sequences.}
Zeno behavior may produce infinitely many rank changes in finite time.
Compactness and regularity hypotheses are needed to define limiting transport.

\paragraph{Hybrid geodesics.}
Existence and uniqueness depend on transition costs, boundary geometry, and
reset rules. No general theorem is claimed here.

\paragraph{Experimental identification.}
Applications require estimators for effective rank, transition maps, energy
defects, and order-dependent loop operators from finite noisy data.

\paragraph{Toroidal closure.}
Boundedness and recurrence do not imply toroidal topology. A torus arises only
under additional hypotheses, such as independent commuting complete phase
flows with a full-rank period lattice. This is logically separate from the
rank-transition holonomy theory.

% ==============================================================================
\section{Conclusion}

Directional Relational Manifolds provide a stratified geometry for systems
whose effective relational dimension changes with state. The foundational
data are
\[
\mathfrak D=
\bigl(M,E,g,\rho,\{S_\alpha\},\{\nabla^\alpha\},\{J_e\}\bigr),
\]
where $E$ is a fixed-rank ambient bundle, $g$ is a smooth
positive-semidefinite metric, $\rho$ is an anchor, $S_\alpha$ are
constant-rank strata, $\nabla^\alpha$ are stratumwise metric connections, and
$J_e$ are rank-transition maps.

The effective fiber and dimension are
\[
\overline E_p=E_p/\Ker g_p,
\qquad
d_{\mathrm{DRM}}(p)=\rank g_p.
\]
Inside a stratum, ordinary metric vector-bundle geometry applies. Between
rank strata, explicit transition events replace a nonexistent globally smooth
fixed-rank connection.

For a closed hybrid loop,
\[
H_\gamma=P_mJ_{m-1}P_{m-1}\cdots J_1P_1
\]
acts on the initial effective fiber. If one transition has rank below the
initial dimension, then $H_\gamma$ is rank-deficient, non-invertible, and
cannot equal the identity. The system may therefore return to the same base
point and the same nominal dimension without recovering its initial
relational state. This is geometric hysteresis induced by dimensional change.

The phenomenon can survive when every regular stratum is flat, showing that
rank history is a source of path dependence distinct from ordinary curvature.
Under dissipativity, the loop operators form a contraction semigroup and admit
rank, norm, energy, spectral, and order-sensitive memory invariants.

Under constant-rank and compatibility hypotheses, DRM reduces to established
Riemannian, sub-Riemannian, and Fisher--Rao geometries. Its proposed
contribution is the integrated treatment of variable metric rank, quotient
geometry, anchored motion, hybrid transition maps, and non-invertible global
transport.

% ==============================================================================
\appendix

\section{Local Matrix Criteria}

In a local frame $e_1,\ldots,e_N$, let
\[
G(p)=\bigl(g_p(e_i,e_j)\bigr).
\]
Then
\[
d_{\mathrm{DRM}}(p)=\rank G(p),
\qquad
K_p=\{z^ie_i:G(p)z=0\}.
\]
On a rank-$r$ stratum, an adapted frame gives
\[
G(p)=\begin{pmatrix}G_{\mathrm{eff}}(p)&0\\0&0\end{pmatrix},
\]
where $G_{\mathrm{eff}}(p)$ is positive definite.

For a transition represented by $J$, dissipativity and isometry are
respectively equivalent to
\[
J^\top G_+J\preceq G_-,
\qquad
J^\top G_+J=G_-.
\]

For a closed loop with initial metric $G$, the energy-defect matrix is
\[
D_\gamma=G-H_\gamma^\top G H_\gamma.
\]
The loop is dissipative exactly when $D_\gamma\succeq0$.

\section{Numerical Diagnostics}

For finite precision, one may estimate effective rank by
\[
d_\varepsilon(G)=\#\{\lambda_i(G)>\varepsilon\}.
\]
A spectral effective rank is
\[
r_{\mathrm{eff}}(G)
=\exp\!\left(-\sum_i p_i\log p_i\right),
\qquad
p_i=\frac{\lambda_i}{\sum_j\lambda_j},
\]
for nonnegative eigenvalues with nonzero sum.

When $G_-$ is positive definite on the effective incoming fiber, define
\[
\delta_J
=\lambda_{\max}\!\left(
G_-^{-1/2}J^\top G_+JG_-^{-1/2}
\right)-1.
\]
Then $\delta_J\le0$ characterizes dissipativity.

Useful loop indices include
\[
h_{\mathrm{rank}}=1-\frac{\rank H_\gamma}{n},
\qquad
h_F=\frac{\|H_\gamma-\Id\|_F}{\sqrt n},
\]
and, for $v\ne0$,
\[
h_{\mathrm{energy}}(v)
=1-\frac{\|H_\gamma v\|^2}{\|v\|^2}.
\]

\section{Conditional Toroidal Closure}

The rank-transition framework does not imply toroidal topology. For
completeness, a sufficient condition is recorded.

\begin{theorem}[Conditional toroidal closure]
Let $A$ be a compact connected invariant manifold admitting $n$ complete,
pointwise linearly independent, commuting vector fields
$X_1,\ldots,X_n$. Suppose their joint $\R^n$-action is transitive and the
stabilizer of a point is a full-rank lattice $\Lambda\subset\R^n$. Then
\[
A\cong\R^n/\Lambda\cong\mathbb T^n.
\]
\end{theorem}

\begin{proof}
The commuting complete fields generate a transitive action of $\R^n$ on $A$.
Thus $A$ is diffeomorphic to the homogeneous quotient
$\R^n/\Stab(p)=\R^n/\Lambda$. A quotient of $\R^n$ by a full-rank lattice is
diffeomorphic to the $n$-torus.
\end{proof}

% ==============================================================================
\section*{Acknowledgments}

Generative artificial intelligence tools were used during the preparation of
this manuscript to assist with drafting, language refinement, structural
review, mathematical consistency checks, and \LaTeX{} editing. The author
reviewed and revised the resulting material and remains fully responsible for
the content, arguments, and conclusions of the work.

\section*{Funding}

This research received no external funding.

\section*{Data Availability}

No data were generated or analyzed for this theoretical work.

\section*{Conflict of Interest}

The author declares no conflict of interest.

\begin{thebibliography}{99}

\bibitem{AmariNagaoka2000}
S.~Amari and H.~Nagaoka,
\emph{Methods of Information Geometry},
American Mathematical Society, 2000.

\bibitem{AndroulidakisSkandalis2009}
I.~Androulidakis and G.~Skandalis,
The holonomy groupoid of a singular foliation,
\emph{Journal f\"ur die reine und angewandte Mathematik}
626 (2009), 1--37.

\bibitem{Chentsov1982}
N.~N. Chentsov,
\emph{Statistical Decision Rules and Optimal Inference},
American Mathematical Society, 1982.

\bibitem{GoebelSanfeliceTeel2012}
R.~Goebel, R.~G. Sanfelice, and A.~R. Teel,
\emph{Hybrid Dynamical Systems: Modeling, Stability, and Robustness},
Princeton University Press, 2012.

\bibitem{KobayashiNomizu1963}
S.~Kobayashi and K.~Nomizu,
\emph{Foundations of Differential Geometry, Volume I},
Wiley, 1963.

\bibitem{Lee2013}
J.~M. Lee,
\emph{Introduction to Smooth Manifolds},
2nd ed., Springer, 2013.

\bibitem{Lee2018}
J.~M. Lee,
\emph{Introduction to Riemannian Manifolds},
2nd ed., Springer, 2018.

\bibitem{Mackenzie2005}
K.~C.~H. Mackenzie,
\emph{General Theory of Lie Groupoids and Lie Algebroids},
Cambridge University Press, 2005.

\bibitem{Nielsen2020}
F.~Nielsen,
An elementary introduction to information geometry,
\emph{Entropy} 22 (2020), 1100.

\bibitem{Stoica2014}
O.~C. Stoica,
On singular semi-Riemannian manifolds,
\emph{International Journal of Geometric Methods in Modern Physics}
11 (2014), 1450041.

\bibitem{Sussmann1973}
H.~J. Sussmann,
Orbits of families of vector fields and integrability of distributions,
\emph{Transactions of the American Mathematical Society}
180 (1973), 171--188.

\end{thebibliography}

\end{document}
